\section{The complementary weight family and primitive dual tests}
\label{sec:cloud}

\status{This section expands the source's Type-II \texttt{CloudStep}: we prove the complementary dimension comparison, the reduction modulo the uniformizer for every primitive combination, the contraction argument, and the valuation bookkeeping. The integral standard-monomial and Weyl-lattice facts, and the constrained signed-chain Gram estimate stated below, remain inputs. A separate source note proves the characteristic-zero operator bound used in the contraction.}

\subsection{A complementary family of the right size}

Use the edge lattices and perfect pairings of Section~\ref{sec:edge-lattices}. Write $k=m+1$, $q=p^{n'}$, $N=|\mathcal T|$, and
\[
 d_s=\begin{cases}\lambda^0_{s,1}-\lambda^0_{s,2},&s\in\mathcal T,\\0,&s\notin\mathcal T,\end{cases}
 \qquad \lambda^{\rm rem}_s=\lambda^0_s-(d_s,0,\ldots,0).
\]
Thus the first two remainder weights agree at every selected twist. Let $D_j(\mu)$ denote the Weyl dimension
$\prod_{i<l}(\mu_i-\mu_l+l-i)/(l-i)$ for a decreasing $j$-tuple. The complementary family has the exact cardinality
\[
 n_2=\prod_{s\in\mathcal T}D_{k-1}(\lambda^0_{s,2},\ldots,\lambda^0_{s,k})
       \prod_{s\notin\mathcal T}D_k(\lambda^0_s).
 \tag{CL.1}\label{eq:cl-count}
\]
Here is the construction behind this formula. Use the integral monic standard-monomial basis of the remainder representation, indexed by interlaced exponent arrays $e(a,t)$, with column sums $\nu_t$; arrays use indices $0,\ldots,k-1$. At a selected even twist retain those with $e(0,0)=\nu_0$. At a selected odd twist retain those with $e(0,t)=0$ for every $t$. At unselected twists retain the entire basis. In the even case deleting the first row and column identifies the selected arrays with the $(k-1)$-variable arrays of tail shape. In the odd case delete the first row and last column; the last column has sum zero after the prescribed dual reversal. The latter shape is the dual of the tail, up to a determinant twist, so it has the same dimension. Tensor these families and transport them to $R_e(\lambda^{\rm rem})$. This gives $h_1,\ldots,h_{n_2}$. \textsf{[CL1]}

Suppose $d_s\ge a>0$ and the spread of the selected tail is at most $c_0$, with $kc_0/a\le1$. Weyl's formula, after subtracting the last weight, gives
\[
 D_k(d+\mu_0,\mu)
 =D_{k-1}(\mu)\prod_{j=1}^{k-1}\frac{d+\mu_0-\mu_{j-1}+j}{j}
 \le D_{k-1}(\mu)(1+c_0/a)^{k-1}\binom{d+k-1}{k-1}.
\]
For $0\le x\le1/k$, $(1-kx)(1+x)^{k-1}\le1$: its logarithmic derivative is negative on $x<1/k$, and the endpoint is immediate. Multiplying the resulting inequalities proves
\[
 (1-kc_0/a)^N\dim V
 \le n_2\prod_{s\in\mathcal T}\binom{d_s+k-1}{k-1}.
 \tag{CL.2}\label{eq:cl-weyl}
\]
Consequently, if the symbol construction supplies a saturated free family $B_1,\ldots,B_b$ of size
$b\ge(1-C_rNr^2k/Y)q^{-r(k-r)}\prod_s\binom{d_s+k-1}{k-1}$, then
\[
 bn_2\ge\frac{\rho\dim V}{q^{r(k-r)}},\qquad
 \rho=(1-C_rNr^2k/Y)(1-kc_0/a)^N.
 \tag{CL.3}
\]
If the first factor is negative this last bound is automatic; otherwise multiply (CL.2) by that factor. The symbol dimension estimate itself is a separate input. \textsf{[CL2]}

\subsection{Why every primitive combination survives}

Let $P$ be the integral polynomial obtained from the source's determinant factor, with $\bar P\ne0$ modulo $w$. Suppose the flat lifts satisfy $\operatorname{subst}(F_b)=B_bP$ in the appropriate renamed variables, and put $S_{b,i}=\Theta(F_b)h_i$. Then
\[
 \exists(b,i):c_{b,i}\in\mathcal O^\times
 \quad\Longrightarrow\quad
 U:=\sum_{b,i}c_{b,i}S_{b,i}\notin wR_e(\lambda^0).
 \tag{CL.4}\label{eq:cl-primitive}
\]
To prove this, first observe why the polynomial products used above are independent over the residue field. At an even twist the flat generators are $X_{c,0}$; the selected remainder's first column is fixed, so the first column of a product's leading exponent recovers every flat exponent $\alpha_c$. At an odd twist they are the signed maximal minors omitting row $c$. The diagonal entries of their leading exponent are $\sum_{c>t}\alpha_c$, $0\le t<k-1$. The selected remainder has zero diagonal: its zero first row and interlacing force $e(a,t)=0$ for $a\le t$. These tail sums, together with $\sum_c\alpha_c=d_s$, recover $\alpha$. The remainder exponent is then recovered by subtraction. Thus all leading monomials are distinct; their leading coefficients are units. Tensoring over the twists preserves independence.

The coordinate change defining $R_e$ sends its elements to integral remainder polynomials, and sends the displayed products to these flat-generator products. Its scalar factors multiply because both the determinant weight and $A$ are additive in the weight. If $U\in wR_e$, reduction modulo $w$ therefore gives
\[
 \sum_i\Theta_0\!\left(\bar P\sum_b\bar c_{b,i}\bar B_b\right)\bar h_i=0.
\]
The independence just proved makes each polynomial in parentheses zero. Cancel $\bar P$ in the residue polynomial ring. Saturation of $\operatorname{span}_{\mathcal O}(B_b)$ and linear independence of the $B_b$ imply $c_{b,i}\in w\mathcal O$ for every pair, a contradiction. This is a test of arbitrary primitive combinations, not just of individual family members. \textsf{[CL3]}

\subsection{Contraction reduces the dual test to a flat representation}

First work on the selected twists alone. The remaining integral representation input supplies maps $\Psi_0,\Psi_f$ from dual to opposite-edge representations satisfying
\[
 p^{\kappa_i}\bar R_i\subseteq\Psi_i(R'_i)\subseteq\bar R_i,
 \qquad i=0,f,
 \tag{CL.5}\label{eq:cl-sandwich}
\]
and their analogous vertex-lattice comparisons. Cartan multiplication spans $\bar R_0$ by $y\beta$, with $y\in\bar R_f$, $\beta\in\bar R_{\rm rem}$. In the standard-monomial model this spanning follows by splitting each product of flag minors into the two prescribed shapes. Set $\kappa=\kappa_0+\kappa_f$ and
$\Lambda_\beta=p^{\kappa_0}\Psi_0^{-1}M_\beta\Psi_f$, where $M_\beta(y)=y\beta$. Define $\rho_\beta$ by
\[
 \beta_f(\rho_\beta u,T)=\beta_0(u,\Lambda_\beta T).
 \tag{CL.6}
\]
Then $\rho_\beta R_0\subseteq R_f$. If all $\rho_\beta u$ lie in $w^{2\kappa+1}R_f$, take any $Z\in R'_0$, express $\Psi_0Z=\sum y\beta$, and use $T=\Psi_f^{-1}(p^{\kappa_f}y)\in R'_f$. Equation (CL.6) makes $p^\kappa\beta_0(u,\Psi_0^{-1}(y\beta))$ divisible by $w^{2\kappa+1}$. Since $p^\kappa=w^{2\kappa}$, summing gives $\beta_0(u,Z)\in w\mathcal O$. Perfectness implies $u\in wR_0$. \textsf{[CL5]}

Contracting the nonselected tensor factors against an integral dual basis preserves this detection statement: if every contraction were divisible by $w$, the original tensor would be divisible by $w$. For the primitive $U$ in (CL.4), therefore choose a contraction and write
$\rho_\beta U=w^\nu\Theta(F_1)$, with $0\le\nu\le2\kappa$ and $\Theta(F_1)$ primitive in $R_f$. The Fock contraction identity places $F_1$ in the span of words in $X_{s,i}\partial_{s,j}$ applied to the $F_b$, using at most $c_0$ letters at each selected twist. A characteristic-zero proof of this assertion, including its highest-weight argument, is given in the expanded note. It makes no assertion that these operator coefficients are integral. \textsf{[CL4]}

Put $\ell=rNc_0+2rn'c''$. The same lattice comparisons at the other edges and at the target give a dual map $\Lambda$ with loss at most $\ell$ in $w$-exponents and a scalar $c\in\mathcal O$ such that
\[
 c\,\beta_0(U,\Lambda T)=p^\kappa\beta_f(\Theta F_1,T).
 \tag{CL.7}\label{eq:cl-transfer}
\]
Explicitly, $\Lambda=w^{2(A_f-A_0)}\rho_\beta^*$ and $c=w^{2\kappa-\nu}$; hence $c$ need not be a unit. The source's nonselected tensor factors can be contracted by an integral dual basis without adding to $\kappa$; they contribute the $2rn'c''$ height loss. For the selected factors,
\[
 \kappa_i=\sum_{s\in\mathcal T,\,h_{i,s}>0}
 \binom{k}{2}\lfloor\log_p(h_{i,s}+k-1)\rfloor,
 \qquad h_{i,s}+k\le a^2,
 \quad\kappa\le4Nk^2\log_p a.
 \tag{CL.8}
\]
Here $h_{i,s}$ is the spread of the corresponding top or flat weight. These are actual lattice comparison exponents, not dimensions of the representations. \textsf{[CL5--CL6]}

\subsection{The flat valuation bound and what remains}

Write $L=Nk^2\log_p a\ge1$ and normalize $v(p)=1$, so $v(w)=1/2$. For a flat polynomial of multidegree $d$, put
\[
 \mathfrak f=\prod_s d_s!,\qquad
 \kappa_{\rm fl}=\sum_s k\lfloor\log_p d_s\rfloor,
 \quad \lfloor\log_p0\rfloor:=0.
\]
The normalized Fock pairing $p^{\kappa_{\rm fl}}\langle G,H\rangle/\mathfrak f$ is integral for integral $G,H$: Legendre's formula bounds the valuation of each relevant multinomial coefficient by $\kappa_{\rm fl}$. The constrained signed-chain construction, still an input here, supplies a test polynomial $H$ with
\[
 \langle G,H\rangle=M+E,\qquad
 v(M)\le v(\mathfrak f)+29L,\qquad
 v(E)\ge c_P-20L+v(\mathfrak f)-\kappa_{\rm fl},
 \tag{CL.9}\label{eq:cl-flat-error}
\]
where $\kappa_{\rm fl}\le20L$ and $c_P>69L$. The strict inequality forces $v(E)>v(M)$, so
\[
 y=p^{\kappa_{\rm fl}}\langle G,H\rangle/\mathfrak f\in\mathcal O,
 \qquad y\mid p^{D_X},\qquad D_X=\lceil49L\rceil\le50L.
 \tag{CL.10}
\]
The expanded note derives (CL.9) from the exact selection, Gram and remainder bounds, and proves the factorial and projection-cost estimates; it explicitly identifies the unexpanded construction of the constrained test. \textsf{[CL7--CL9]}

The localizer realization writes $H$ as a sum of allowed symbols and realizes its pairing as $y+w^{c_P}z$. If $yb=p^{D_X}$ and $2D_X<c_P$, then
\[
 (y+w^{c_P}z)b=p^{D_X}\bigl(1+w^{c_P-2D_X}zb\bigr);
\]
the parenthesized factor is a unit, so the realized flat value still divides $p^{D_X}$. Applying (CL.7), an integral full value $x=\beta_0(U,\Lambda T)$ satisfies $cx=p^\kappa y'$, where $y'b'=p^{D_X}$. Thus $x(cb')=p^{\kappa+D_X}$. We obtain, for \emph{every} primitive coefficient array,
\[
 D_1=\kappa+D_X\le54Nk^2\log_p a,
 \qquad t_T\le t_0+Nc_0+\ell,
 \qquad \Lambda T\in w^{c_*}R'_{e'}\quad(e'\ne e),
 \tag{CL.11}\label{eq:cl-final}
\]
using $c_P=c_*+\ell$. The source hypotheses include $D_1\le D_p$ and $c_*\ge8(D_p+1)$ once its large constant exceeds the displayed bounds. They ensure the needed strict perturbation inequalities. This extracts the usable constant $54$ from the pinned witness choices; it does not remove the remaining integral Weyl comparison, symbol dimension, or constrained signed-chain Gram input. \textsf{[CL10--CL12]}
